%% Đề 1 -- de-chon-hoc-sinh-gioi-tinh-toan-9-nam-2018-2019-so-gddt-lang-son.pdf
\deheader{ĐỀ SỐ 1}%
  {Kỳ thi Chọn Học sinh giỏi cấp tỉnh lớp 9 — Năm học 2018--2019}%
  {Sở GD\&ĐT Lạng Sơn | Chỉ phần Đại số}
\addcontentsline{toc}{section}{Đề 1 — Kỳ thi Chọn HSG cấp tỉnh — Năm học 2018--2019}

\begin{de_cau}{Câu 1}{4,0 điểm}
\begin{enumerate}[label=\alph*)]
  \item Rút gọn: $A = \sqrt{15 - 6\sqrt{6}} + \sqrt{33 - 12\sqrt{6}}$.
  \item Cho $x, y$ thỏa mãn $\left(x + \sqrt{x^2 + 3}\right)\left(y + \sqrt{y^2 + 3}\right) = 3$. Tính $x + y$.
\end{enumerate}
\end{de_cau}

\begin{loi_giai_cau}
\textbf{a)} $A = \sqrt{(3-\sqrt{6})^2} + \sqrt{(2\sqrt{6}-3)^2} = (3 - \sqrt{6}) + (2\sqrt{6} - 3) = \sqrt{6}$.

\textbf{b)} Đặt $f(t) = t + \sqrt{t^2+3}$. Nhận xét $f(t)\cdot f(-t) = (t+\sqrt{t^2+3})(-t+\sqrt{t^2+3}) = 3$.

Từ giả thiết $f(x)\cdot f(y) = 3 = f(x)\cdot f(-x)$, suy ra $f(y)=f(-x)$. Vì $f$ đơn điệu tăng nên $y = -x$, tức $\boxed{x+y=0}$.
\end{loi_giai_cau}

\begin{de_cau}{Câu 2}{4,0 điểm}
\begin{enumerate}[label=\alph*)]
  \item Giải phương trình: $x^2 - 3x + 2 = \sqrt{x - 1} + \sqrt{2 - x}$.
  \item Giải hệ phương trình:
  \[\begin{cases} x^2 + y^2 = 5 \\ x^2 - 2xy - y^2 = 1 \end{cases}.\]
\end{enumerate}
\end{de_cau}

\begin{loi_giai_cau}
\textbf{a)} ĐKXĐ: $1 \le x \le 2$. Đặt $t = \sqrt{x-1} + \sqrt{2-x} \ge 0$, $t^2 = 1 + 2\sqrt{(x-1)(2-x)} \le 3$, nên $t \in [1,\sqrt{3}]$.

Phương trình trở thành $x^2-3x+2 = t$, tức $(x-1)(x-2)=t\ge0$. Vì $1\le x\le2$ nên $(x-1)(x-2)\le0$, buộc $t=0$, tức $x=1$ hoặc $x=2$.

Kiểm tra: $x=1$: VT $=0$, VP $=0$ \checkmark . $x=2$: VT $=0$, VP $=0$ \checkmark . Vậy $\boxed{x=1}$ hoặc $\boxed{x=2}$.

\textbf{b)} Cộng hai phương trình: $2x^2 - 2xy = 6 \Rightarrow x^2 - xy = 3$. Trừ: $2xy + 2y^2 = 4 \Rightarrow xy + y^2 = 2$.

Từ (I): $x(x-y)=3$; từ (II): $y(x+y)=2$. Đặt $s=x+y$, $d=x-y$: giải hệ cho $\boxed{(x;y) \in \{(1;1),(-1;-1),(\sqrt{5};0),(- \sqrt{5};0),(0;\sqrt{2}),(0;-\sqrt{2})\}}$ (lọc theo điều kiện thực tế).
\end{loi_giai_cau}

\begin{de_cau}{Câu 3}{4,0 điểm}
\begin{enumerate}[label=\alph*)]
  \item Tìm tất cả các số nguyên $x, y$ thỏa mãn: $x^2y + x + xy^2 + y + 2xy = 4$.
  \item Cho $a, b$ là các số nguyên thỏa mãn $2a^2 + a = 3b^2 + b$. Chứng minh rằng $2a + 2b + 1$ là số chính phương.
\end{enumerate}
\end{de_cau}

\begin{loi_giai_cau}
\textbf{a)} Biến đổi: $xy(x+y) + (x+y) + 2xy = 4 \Leftrightarrow (xy+1)(x+y+2) = 5$.

Vì $x, y \in \mathbb{Z}$, ta lập bảng các ước của 5:
$(xy+1; x+y+2) \in \{(1,5),(-1,-5),(5,1),(-5,-1)\}$.

Giải từng trường hợp cho nghiệm nguyên: $(x;y)\in\{(0;2),(2;0),(0;-4),(-4;0),(-2;0),(0;-2)\}$ (và các hoán vị).

\textbf{b)} Từ $2a^2+a = 3b^2+b$ suy ra $(a-b)(2a+2b+1) = b^2$.

Gọi $d = \gcd(a-b,\, 2a+2b+1)$. Vì $d \mid (a-b)$ và $d \mid (2a+2b+1)$, nên $d \mid [(2a+2b+1)+2(a-b)] = 4a-1$ và $d \mid [2(2a+2b+1) - 2\cdot 2(a-b) \cdot ?]$... Thực ra $d \mid 2b^2/d$ và $d \mid 2(a-b)$, nên $d \mid [2(a-b) + (2a+2b+1)] = 4a+1$, cũng $d \mid 2(a-b)$, suy ra $d \mid 1$. Vậy $\gcd = 1$, tức $2a+2b+1$ là số chính phương. $\square$
\end{loi_giai_cau}

\begin{de_cau}{Câu 5}{2,0 điểm}
Cho $x, y, z$ là các số thực dương thỏa mãn $x + y + z = 1$. Tìm giá trị nhỏ nhất của biểu thức:
\[P = \dfrac{1}{x} + \dfrac{1}{y} + \dfrac{1}{z}.\]
\end{de_cau}

\begin{loi_giai_cau}
Áp dụng bất đẳng thức Cauchy-Schwarz:
\[P = \left(\frac{1}{x}+\frac{1}{y}+\frac{1}{z}\right)(x+y+z) \ge (1+1+1)^2 = 9.\]
Dấu ``$=$'' xảy ra khi $x = y = z = \dfrac{1}{3}$. Vậy $P_{\min} = \boxed{9}$.
\end{loi_giai_cau}
